% geometry/curvature_classification.tex

\chapter{Curvature Classification}

\label{chap:curvature-classification}

% ============================================================
\section*{Motivation}
% ============================================================

Classical railway alignment theory typically classifies transitions by
their analytical formulas or historical engineering usage.

Within AIM, transition and alignment structures are classified
intrinsically within curvature space.

This allows classification based on:
\begin{itemize}
\item continuity,
\item curvature behaviour,
\item dynamics,
\item realization behaviour,
\item spectral structure,
\item runtime robustness,
\item and operational suitability.
\end{itemize}

Curvature classification therefore becomes a structural interpretation
framework rather than merely a catalogue of curve formulas.

\begin{principle}[Classification as interpretation]
\label{princ:classification-as-interpretation}
Within AIM, classification is not a naming scheme for curve formulas.
It is an interpretation layer for curvature behaviour.
\end{principle}

% ============================================================
\section{Classification Principle}
% ============================================================

\begin{principle}[Intrinsic classification principle]
\label{princ:intrinsic-classification}
Within AIM, alignments are classified by intrinsic curvature behaviour
rather than by reconstructed position geometry.
\end{principle}

Two geometrically similar curves in position space may possess very
different curvature behaviour.

Conversely, geometrically different realizations may belong to the same
curvature class if their intrinsic curvature evolution behaves
equivalently.

Classification is therefore performed in curvature space before it is
interpreted in position space.

% ============================================================
\section{Continuity Classification}
% ============================================================

\subsection{\texorpdfstring{$G^0$}{G0} Geometry}

\begin{definition}[\(G^0\)-continuous geometry]
\label{def:g0-geometry}
A geometry is \(G^0\)-continuous if position is continuous.
\end{definition}

% ------------------------------------------------------------

\subsection{\texorpdfstring{$G^1$}{G1} Geometry}

\begin{definition}[\(G^1\)-continuous geometry]
\label{def:g1-geometry}
A geometry is \(G^1\)-continuous if tangent direction is continuous.
\end{definition}

% ------------------------------------------------------------

\subsection{\texorpdfstring{$G^2$}{G2} Geometry}

\begin{definition}[\(G^2\)-continuous geometry]
\label{def:g2-geometry}
A geometry is \(G^2\)-continuous if curvature is continuous.
\end{definition}

Within railway alignment modelling, \(G^2\)-continuity is often the
minimal requirement for dynamically acceptable transitions.

% ------------------------------------------------------------

\subsection{\texorpdfstring{$G^3$}{G3} Geometry}

\begin{definition}[\(G^3\)-continuous geometry]
\label{def:g3-geometry}
A geometry is \(G^3\)-continuous if curvature gradient is continuous.
\end{definition}

Higher continuity classes improve:
\begin{itemize}
\item jerk behaviour,
\item realization smoothness,
\item runtime stability,
\item and vehicle response.
\end{itemize}

Continuity classification bridges classical geometry notation with
curvature-space interpretation.

% ============================================================
\section{Curvature Behaviour Classification}
% ============================================================

\subsection{Constant Curvature Class}

\begin{definition}[Constant curvature class]
\label{def:constant-curvature-classification}
A curvature law belongs to the constant class if:
\[
\kappa(s)=\kappa_0.
\]
\end{definition}

% ------------------------------------------------------------

\subsection{Linear Curvature Class}

\begin{definition}[Linear curvature class]
\label{def:linear-curvature-class}
A curvature law belongs to the linear class if:
\[
\kappa(s)=a s+b.
\]
\end{definition}

The classical clothoid belongs to this class.

% ------------------------------------------------------------

\subsection{Polynomial Curvature Class}

\begin{definition}[Polynomial curvature class]
\label{def:polynomial-curvature-class}
A curvature law belongs to the polynomial class if:
\[
\kappa(s)
=
\sum_i a_i s^i.
\]
\end{definition}

% ------------------------------------------------------------

\subsection{Trigonometric Curvature Class}

\begin{definition}[Trigonometric curvature class]
\label{def:trigonometric-curvature-class}
A curvature law belongs to the trigonometric class if it contains
sinusoidal or cosine-based components.
\end{definition}

% ------------------------------------------------------------

\subsection{Lookup-Based Curvature Class}

\begin{definition}[Lookup-based curvature class]
\label{def:lookup-curvature-class}
A curvature law belongs to the lookup-based class if it is defined by
sampled or tabulated curvature values together with interpolation and
reconstruction rules.
\end{definition}

% ------------------------------------------------------------

\subsection{Hybrid Curvature Class}

\begin{definition}[Hybrid curvature class]
\label{def:hybrid-curvature-class}
A hybrid curvature law combines:
\begin{itemize}
\item parametric structure,
\item lookup structure,
\item measured data,
\item and sampled corrections.
\end{itemize}
\end{definition}

% ============================================================
\section{Symmetry Classification}
% ============================================================

\subsection{Symmetric Transitions}

\begin{definition}[Symmetric transition]
\label{def:classification-symmetric-transition}
A normalized transition is symmetric if its curvature evolution is
symmetric with respect to the midpoint of the normalized interval.
\end{definition}

For transition comparison, symmetry should usually be evaluated in
normalized transition space rather than in raw physical coordinates.

% ------------------------------------------------------------

\subsection{Antisymmetric Structure}

\begin{definition}[Antisymmetric transition]
\label{def:classification-antisymmetric-transition}
A transition has antisymmetric structure if its deviation from a
reference or midpoint value changes sign under midpoint reflection.
\end{definition}

Antisymmetry is especially useful when analysing centred transition
shape functions or derivative behaviour.

% ------------------------------------------------------------

\subsection{Asymmetric Transitions}

\begin{definition}[Asymmetric transition]
\label{def:classification-asymmetric-transition}
A transition is asymmetric if no symmetry condition is imposed.
\end{definition}

Asymmetry may arise intentionally through:
\begin{itemize}
\item realization-aware design,
\item vehicle-dynamics optimization,
\item topology constraints,
\item or local engineering constraints.
\end{itemize}

% ============================================================
\section{Derivative Classification}
% ============================================================

\subsection{Bounded Gradient Class}

\begin{definition}[Bounded gradient class]
\label{def:bounded-gradient-class}
A curvature law belongs to the bounded-gradient class if:
\[
|\kappa'(s)|
\le
\eta_{\max}.
\]
\end{definition}

% ------------------------------------------------------------

\subsection{Bounded Higher-Derivative Class}

\begin{definition}[Bounded higher-derivative class]
\label{def:bounded-higher-derivative-class}
A curvature law belongs to this class if:
\[
|\kappa''(s)|<\infty,
\qquad
|\kappa'''(s)|<\infty.
\]
\end{definition}

These classes influence:
\begin{itemize}
\item dynamic smoothness,
\item numerical stability,
\item realization robustness,
\item and optimization conditioning.
\end{itemize}

% ============================================================
\section{Spectral Classification}
% ============================================================

\subsection{Low-Frequency Structures}

\begin{definition}[Low-frequency curvature class]
\label{def:low-frequency-curvature-class}
Low-frequency curvature structures are dominated by slowly varying
curvature evolution.
\end{definition}

Low-frequency structure usually determines the large-scale geometric
shape of the alignment.

% ------------------------------------------------------------

\subsection{High-Frequency Structures}

\begin{definition}[High-frequency curvature class]
\label{def:high-frequency-curvature-class}
High-frequency curvature structures contain rapidly varying or
oscillating components.
\end{definition}

High-frequency structures tend to:
\begin{itemize}
\item amplify realization sensitivity,
\item increase dynamic roughness,
\item increase measurement sensitivity,
\item and reduce numerical stability.
\end{itemize}

\begin{proposition}[Spectral realization principle]
\label{prop:classification-spectral-realization}
Realization suppresses high-frequency curvature components.
\end{proposition}

Spectral classification therefore links transition theory with physical
realization behaviour.

% ============================================================
\section{Realization Classification}
% ============================================================

\subsection{Realization-Stable Class}

\begin{definition}[Realization-stable class]
\label{def:realization-stable-class}
A transition belongs to the realization-stable class if:
\begin{itemize}
\item dominant curvature structure is preserved under realization,
\item low-frequency behaviour remains identifiable,
\item and realization does not introduce instability.
\end{itemize}
\end{definition}

% ------------------------------------------------------------

\subsection{Realization-Sensitive Class}

\begin{definition}[Realization-sensitive class]
\label{def:realization-sensitive-class}
A transition belongs to the realization-sensitive class if small
realization perturbations strongly alter curvature behaviour.
\end{definition}

Realization-sensitive transitions may be mathematically valid but
operationally undesirable.

% ============================================================
\section{Runtime Classification}
% ============================================================

\subsection{Projection-Stable Class}

\begin{definition}[Projection-stable class]
\label{def:projection-stable-class}
A curvature structure is projection-stable if:
\begin{itemize}
\item world-to-track projection remains numerically robust,
\item local inverse problems remain well-conditioned,
\item and ambiguity remains controllable.
\end{itemize}
\end{definition}

% ------------------------------------------------------------

\subsection{Runtime-Smooth Class}

\begin{definition}[Runtime-smooth class]
\label{def:runtime-smooth-class}
A curvature structure is runtime-smooth if:
\begin{itemize}
\item frame evaluation remains continuous,
\item pose reconstruction remains robust,
\item visualization remains stable,
\item and runtime operators behave predictably.
\end{itemize}
\end{definition}

% ============================================================
\section{Operational Classification}
% ============================================================

\subsection{Comfort-Oriented Class}

\begin{definition}[Comfort-oriented class]
\label{def:comfort-oriented-class}
A transition belongs to the comfort-oriented class if:
\begin{itemize}
\item jerk remains low,
\item lateral acceleration evolves smoothly,
\item and dynamic excitation is minimized.
\end{itemize}
\end{definition}

% ------------------------------------------------------------

\subsection{Realization-Oriented Class}

\begin{definition}[Realization-oriented class]
\label{def:realization-oriented-class}
A transition belongs to the realization-oriented class if:
\begin{itemize}
\item construction robustness is prioritized,
\item machine realization remains stable,
\item maintenance sensitivity is reduced,
\item and physical filtering preserves the intended behaviour.
\end{itemize}
\end{definition}

% ------------------------------------------------------------

\subsection{Optimization-Oriented Class}

\begin{definition}[Optimization-oriented class]
\label{def:optimization-oriented-class}
A transition belongs to the optimization-oriented class if:
\begin{itemize}
\item parameterization is numerically stable,
\item gradients remain well-conditioned,
\item constraints are expressible,
\item and sparse representation is efficient.
\end{itemize}
\end{definition}

% ============================================================
\section{Family Classification}
% ============================================================

\begin{definition}[Transition family classification]
\label{def:transition-family-classification}
Transition families may be classified according to:
\begin{itemize}
\item analytical form,
\item continuity order,
\item derivative behaviour,
\item spectral behaviour,
\item realization behaviour,
\item symmetry,
\item runtime stability,
\item and operational suitability.
\end{itemize}
\end{definition}

Thus, transition classification becomes multidimensional rather than
formula-based.

A single transition family may be favourable in one classification axis
and unfavourable in another.

% ============================================================
\section{Classification Vectors}
% ============================================================

\begin{definition}[Classification vector]
\label{def:classification-vector}
A transition may be assigned a classification vector:
\[
\chi(T)
=
(
\chi_{\mathrm{cont}},
\chi_{\mathrm{deriv}},
\chi_{\mathrm{spec}},
\chi_{\mathrm{real}},
\chi_{\mathrm{run}},
\chi_{\mathrm{op}}
),
\]
where the components describe continuity, derivative behaviour,
spectral structure, realization behaviour, runtime stability, and
operational suitability.
\end{definition}

The classification vector is not intended as a rigid taxonomy. It is a
structured way to compare transition families across multiple relevant
dimensions.

% ============================================================
\section{Canonical Interpretation}
% ============================================================

\begin{proposition}
Within AIM, curvature classification is fundamentally intrinsic rather
than coordinate-based.
\end{proposition}

\begin{proposition}
Transition quality cannot be characterized sufficiently by position
geometry alone.
\end{proposition}

\begin{proposition}
Realization, runtime stability, dynamics, and optimization require
classification directly in curvature space.
\end{proposition}

\begin{proposition}
Transition taxonomy in AIM is multidimensional and context-dependent.
\end{proposition}

% ============================================================
\section{Connection to AIM}
% ============================================================

\begin{summary}
Curvature classification provides the structural interpretation framework
of AIM geometry.

Within this framework:
\begin{itemize}
\item transitions are classified intrinsically,
\item formula names are secondary to curvature behaviour,
\item realization behaviour becomes analyzable,
\item runtime robustness becomes class-dependent,
\item operational suitability becomes comparable,
\item and optimization acts directly on curvature structures.
\end{itemize}

This establishes curvature space as the primary classification domain of
railway alignment geometry and provides the basis for a taxonomy of
transition families.
\end{summary}
